\documentclass[]{ceurart}

%%
%% One can fix some overfulls
\sloppy

%%
%% Minted listings support
%% Need pygment <http://pygments.org/> <http://pypi.python.org/pypi/Pygments>
\usepackage{listings}
%% auto break lines
\lstset{breaklines=true}

%%
%% end of the preamble, start of the body of the document source.
\begin{document}

%%
%% Rights management information.
%% CC-BY is default license.
\copyrightyear{2026}
\copyrightclause{Copyright for this paper by its authors.
  Use permitted under Creative Commons License Attribution 4.0
  International (CC BY 4.0).}

%%
%% This command is for the conference information
\conference{10th Workshop on Practical Aspects of Automated Reasoning,
Lisboa, 2026}

%%
%% The "title" command
\title{Efficient Multi-Scale Indexing with Dynamic IntMaps}


%%
%% The "author" command and its associated commands are used to define
%% the authors and their affiliations.

\author[1]{Albert Eisfeld}[%
email=inf23040@lehre.dhbw-stuttgart.de,
]

\author[1]{Stephan Schulz}[%
orcid=0000-0001-6262-8555,
email=schulz@eprover.org,
]


\address[1]{DHBW Stuttgart, Stuttgart, Germany}
%%
%% The abstract is a short summary of the work to be presented in the
%% article.
\begin{abstract}
  Term and clause indexing are central technologies for efficient
  first-order theorem provers and similar reasoning systems. While
  there are a selection of basic technologies, most of them require
  the efficient mapping from function symbols (usually encoded as
  small integers) or other integer values to branches of a tree
  structure. Today, both the range of possible value and the
  cardinality of each individual map cary wildly. We present a data
  structure that is efficient in time and space at all scales and key
  distributions.
\end{abstract}

%%
%% Keywords. The author(s) should pick words that accurately describe
%% the work being presented. Separate the keywords with commas.
\begin{keywords}
  term indexing \sep
  data structures \sep
  theorem proving \sep
\end{keywords}

%%
%% This command processes the author and affiliation and title
%% information and builds the first part of the formatted document.
\maketitle

\section{Introduction}

Most powerful modern theorem provers like E~\cite{SCV:CADE-2019},
Vampire~\cite{BBCHHHKRRRSSV:CAV-2025} and iProver~\cite{DK:IJCAR-2020}
are based on saturating calculi, in particularly variations of
superposition~\cite{BG94} and related calculi. These calculi represent
the proof state as a set of clauses. They enrich the proof state by
adding logical consequences of existing clauses, with the aim of
eventually generating the empty clause as an explicit witness of
unsatisfiability (and hence conclude a proof by contradiction). In
addition to generating inferences, simplifying inferences like
subsumption and rewriting are critical for success. These techniques
either remove redundant clause, or replace them by simpler ones, and
typically take up most of the time spent by the inference engine.

Successful provers combine efficient execution and intelligent proof
search heuristics to achieve their power. Especially for harder
problems with large search states, \emph{term indexing} is a critical
technology. Indexing of terms (and of clauses) allows a system to
quickly identify potential inference partners for generating
inferences like paramodulation/superposition or resolution, but in
particular also for simplification techniques like rewriting and
subsumption.

Early contributions were made by Stickel~\cite{Stickel:TRIndexing-89},
who introduced \emph{path indexing}, and by McCune, who integrated
path indexing and \emph{discrimination tree indexing} into
Otter~\cite{Mc92}, creating the first successful modern saturating
prover.

Our own prover E uses \emph{perfect discrimination tree
  indexing} for forward rewriting, \emph{fingerprint
  indexing}~\cite{Schulz:IJCAR-2012} for backwards rewriting and
superposition/paramodulation, and \emph{feature vector
  indexing}~\cite{Schulz:FVI-2013} for subsumption and related
techniques.

\cite{ST:JACM-85}

Motivation!

\section{Motivation}


\section{Conclusion}




%% The declaration on generative AI comes in effect
%% in Janary 2025. See also
%% https://ceur-ws.org/GenAI/Policy.html
\section*{Declaration on Generative AI}

The author(s) have not employed any Generative AI tools.

%%
%% Define the bibliography file to be used
\bibliography{stsbib.bib}

%%
%% If your work has an appendix, this is the place to put it.


\end{document}

%%
%% End of file

%%% Local Variables:
%%% mode: LaTeX
%%% TeX-master: t
%%% End:
